\sectionTitle{Experience}{}

\begin{experiences}

 \researchexperience
    {Present}   {Amazon Lab126}{San Francisco, USA}{}
    {April 2025} {\textit{Senior Applied Scientist}\\
     Work on \textbf{extreme hardware-software co-design}, leading Generative AI capabilities on Amazon's custom silicon, a role directly enabled by the acquihire of Nyun AI. Scaled model compression and hardware-aware optimization pipelines developed at Nyun AI to enable on-device GenAI for \href{https://www.aboutamazon.com/news/devices/amazon-new-echo-devices-alexa-plus}{Alexa+}.
    }
  \emptySeparator

 \researchexperience
    {March 2025}   {Nyun AI}{Bangalore, India}{}
    {Aug 2023} {\textit{Co-founder and CTO | Advisor/CEO:  \href{https://dkgupta90.github.io/}{Dr. Deepak Gupta}}\\
     Led the development of advanced compression pipelines for foundational models and optimized their deployment on edge devices while managing a multidisciplinary team to deliver tailored solutions for Fortune 5 enterprises. \textbf{Acquihired by Amazon Lab126.} Published work on model compression and efficiency[\hyperref[C7]{ACL'24}, \hyperref[J3]{TMLR'24, '25}, \hyperref[C6]{ICLR'24}, \hyperref[C7]{IJCAI'24}].
    }
  \emptySeparator


  % \researchexperience
  %   {Jul 2023}   {Mohammed Bin Zayed University of AI}{Abu Dhabi, UAE}{}
  %   {Mar 2023} {\textit{Research Assistant | Advisor: \href{https://zhiqiangshen.com/}{Prof. Zhiqiang Shen}}, \href{https://www.cs.cmu.edu/~epxing/}{Prof. Eric Xing}\\
  %   Designed a novel approach - GLoRA, to enhance parameter-efficient fine-tuning (PEFT) methodologies for large-scale image and language models, focused on improving the generalization capabilities of current PEFT techniques.
  %   }
  % \emptySeparator

      \researchexperience
    {Feb 2023}   {Microsoft Research}{Bangalore, India}{}
    {Aug 2022} {\textit{Research Intern | Advisor: \href{https://www.microsoft.com/en-us/research/people/susitara/}{Dr. Sunayana Sitaram}, \href{https://www.microsoft.com/en-us/research/people/supartha/}{Suresh Parthasarathy}}\\
    Conducted research on large-scale HTML understanding models to improve generalization in information extraction from web content and investigated the implications of model compression on bias and fairness in large language models, published at [\hyperref[C5]{ACL'23}].
    }
  % \emptySeparator
  %     \researchexperience
  %   {Jul 2022}   {Amazon Science}{Bangalore, India}{}
  %   {May 2022} {\textit{Applied Science Intern | Advisor: \href{https://in.linkedin.com/in/prateeksircar}{Prateek Sircar}}\\
  %   Contributed to the Aspect Insights service by developing algorithms for unsupervised extraction and clustering of product aspects from textual data. Engineered a framework for integration with automated product video generation to enhance user engagement through visual content.
  %   }
  \emptySeparator
    \researchexperience
    {Apr 2022}   {Google DeepMind}{Bangalore, India}{}
    {Jan 2022} {\textit{Student Researcher | Advisor: \href{https://sites.google.com/site/pshenoyuw/}{Dr. Pradeep Shenoy}}\\
    Executed a comprehensive empirical study on simplicity bias in deep neural networks, establishing methodologies to quantify and mitigate this bias, and developed synthetic datasets and experimental protocols to enhance the understanding of neural network performance dynamics.
    }
  \emptySeparator

    \researchexperience
    {Dec 2021}   {Mohammed Bin Zayed University of AI}{Abu Dhabi, UAE}{}
    {Aug 2021} {\textit{Research Assistant | Advisor: \href{https://zhiqiangshen.com/}{Prof. Zhiqiang Shen}}\\
    Innovated methods for compressing vision transformers by introducing sparsity across transformer modules within a unified framework, demonstrating significant advancements over existing compression techniques. Results published at [\hyperref[C3]{CVPR'22}. \hyperref[S1]{TMLR'25}].
    }
  % \emptySeparator
  % \researchexperience
  %   {Jul 2021}   {Amazon Science}{Bangalore, India}{}
  %   {May 2021} {\textit{Applied Science Intern | Advisor: \href{https://www.linkedin.com/in/saurabhshekhar/}{Saurabh Shekhar}}\\
  %   Developed a novel metric, Trivial Modification Rate (TMR), to quantify trivial alterations in product search queries, enhancing the efficacy of search algorithms by 15\%. TMR implementation in production environments facilitated cross-team evaluations of operational impacts on search performance.
  %   }
    
  % \emptySeparator
  %     \researchexperience
  %   {Feb 2021}   {RWTH Aachen University}{Aachen, Germany}{}
  %   {Dec 2020} {\textit{Research Assistant | Advisor: \href{https://www.lfb.rwth-aachen.de/en/institute/ehemalige-lehrstuhlinhaberinnen/merhof/}{Prof. Dorit Mehrof}}\\
  %   Advanced unsupervised anomaly detection methodologies by fine-tuning learned representations through the induction of multivariate Gaussian distributions. Results published at [\hyperref[C1]{IMPROVE'21}].
  %   }
  % \emptySeparator
  %   \researchexperience
  %   {Oct 2020}   {Bosch Research}{Bangalore, India}{}
  %   {Aug 2020} {\textit{Research Intern | Advisor: \href{https://in.linkedin.com/in/dr-sahana-prabhu-295b3382}{Dr. Sahana Prabhu}}\\
  %   Developed a comprehensive framework for semi-supervised and unsupervised anomaly detection and localization across diverse datasets, contributing to enhancements of existing models for ground terrain datasets in outdoor scene analysis.
  %   }
      % \emptySeparator


  %     \researchexperience
  %   {Jun 2020}   {Univeristy of Waterloo}{Waterloo, Canada}{}
  %   {Apr 2020} {\textit{Research Intern | Advisor: \href{https://sites.google.com/view/teellabvengulab/home}{Prof. Vasudevan Lakshminarayanan}}\\
  %   Proposed a method for quantifying and extracting Retinal Blood Vessel Characteristics (RBVC) from Optical Coherence Tomography Angiography (OCTA) images, with research findings accepted at [Photonics West 2021, SPIE].
  %   }
  \emptySeparator
  \researchexperience
    {May 2023}   {Transmute AI Research}{Tromsø, Norway}{}
    {May 2020} {\textit{Co-founding member, Senior Researcher | Advisor: \href{https://dkgupta90.github.io/}{Dr. Deepak K. Gupta}, \href{https://sites.google.com/site/dilipprasad/}{Prof. Dilip K. Prasad}}\\
    \textbf{Co-founded} a research lab to foster research culture among peers and aspiring researchers, solicited funding by \href{https://texmin.in/}{Texmin Hub} and \href{https://www.bioailab.org/}{Bio-AI Lab} of UiT Norway. \newline
    Guided \textbf{10+ UG students} from different colleges to pursue research. \newline
    Worked on research problems around network compression, meta-learning and efficient ML leading to 4 publications so far [\hyperref[C1]{ICLR'21}, \hyperref[C2]{ICIP'21}, \hyperref[C4]{CVPR'22}, ICASSP'23].
    }
  \emptySeparator
          
\end{experiences}
\vspace{-2mm}
